\documentclass[10pt,a4paper]{article}
\usepackage[utf8]{inputenc}
\usepackage[T1]{fontenc}
\usepackage{amsfonts}
\usepackage{enumitem}
\usepackage{amssymb}
\usepackage{amsmath}
\usepackage[left=1cm,right=1cm,top=1.5cm,bottom=1.5cm]{geometry}
\usepackage{multirow}
\usepackage[dvipsnames]{xcolor}
\usepackage{graphicx}
\setlength{\parindent}{0pt}
\usepackage{fancyhdr}
\pagestyle{fancy}
\renewcommand\headrulewidth{1pt}
\fancyhead[L]{\small Lycée Denis Diderot }
\fancyhead[R]{\small Classe de Terminale}
\setcounter{page}{1}\fancyfoot[C]{\small page \thepage}
\fancyfoot[R]{Année 2025 2026}
\fancyfoot[L]{\small Alexis RUIZ}
\setlength{\headheight}{14pt}

\begin{document}
\begin{center}
\vspace{-5mm}
\Large{Questionnaire : Tesla et la guerre des courants -- \textcolor{red}{CORRIGÉ}} \\
\vspace{1mm}
\normalsize
\hrule
\end{center}
\vspace{1mm}
{\small A partir de la BD de Sciences et Vie junior N°230 de novembre 2008, répondre aux questions si dessous :}
\vspace{1mm}

\begin{enumerate}[itemsep=1mm]
 
\item Quelle est la nationalité de Nicola Tesla ? \hfill \textcolor{red}{(0,5 pt)}

{\small \textcolor{blue}{\textbf{Réponse :} Nicola Tesla est de nationalité serbe, né en 1856 dans l'Empire austro-hongrois (actuelle Croatie).}}

\item Quel est son métier ? \hfill \textcolor{red}{(0,5 pt)}

{\small \textcolor{blue}{\textbf{Réponse :} Il est ingénieur et inventeur, spécialisé dans le domaine de l'électricité et de l'électromagnétisme.}}

\item Décrire brièvement l'idée géniale de N Tesla ? \hfill \textcolor{red}{(1 pt)}

{\small \textcolor{blue}{\textbf{Réponse :} Son idée géniale est l'invention du courant alternatif et du moteur à induction qui fonctionne grâce à un champ magnétique rotatif, permettant ainsi de transporter l'électricité sur de longues distances de manière efficace.}}

\item Quels sont les inconvénients du courant continu ? \hfill \textcolor{red}{(1 pt)}

{\small \textcolor{blue}{\textbf{Réponse :} Le courant continu présente d'importantes pertes d'énergie lors du transport sur de longues distances. Il nécessite des stations de relais tous les 2 à 3 km pour maintenir la tension, ce qui rend le système coûteux et peu pratique.}}

\item  Quel personnage célèbre part rencontrer N Tesla aux États-Unis ? \hfill \textcolor{red}{(0,5 pt)}

{\small \textcolor{blue}{\textbf{Réponse :} Thomas Edison, le célèbre inventeur américain, embauche Tesla dans son entreprise aux États-Unis.}}

\item Quelle est la différence entre un rotor et un stator ? \hfill \textcolor{red}{(1 pt)}

{\small \textcolor{blue}{\textbf{Réponse :} Le stator est la partie fixe d'un moteur électrique qui génère un champ magnétique, tandis que le rotor est la partie mobile qui tourne à l'intérieur du stator sous l'effet de ce champ magnétique.}}

\item Quelle est la source du conflit entre N Tesla et T Edison ? \hfill \textcolor{red}{(1 pt)}

{\small \textcolor{blue}{\textbf{Réponse :} Le conflit provient du fait qu'Edison refuse de payer la prime promise à Tesla (environ 50 000 dollars) pour l'amélioration de ses dynamos. De plus, Edison est un fervent défenseur du courant continu et s'oppose à l'idée du courant alternatif de Tesla.}}

\item Quel est le nom de l'industriel qui va engager N Tesla comme consultant ? \hfill \textcolor{red}{(0,5 pt)}

{\small \textcolor{blue}{\textbf{Réponse :} George Westinghouse, industriel américain et fondateur de la Westinghouse Electric Company, engage Tesla comme consultant et rachète ses brevets.}}

\item Quel est l'avantage de la tension alternative ? \hfill \textcolor{red}{(1,5 pt)}

{\small \textcolor{blue}{\textbf{Réponse :} L'avantage majeur du courant alternatif est qu'il peut être facilement transformé à l'aide de transformateurs. On peut augmenter la tension pour le transport (réduisant ainsi les pertes en ligne), puis la diminuer pour une utilisation domestique, ce qui permet de transporter l'électricité sur de très longues distances de manière économique et efficace.}}

\item A quelle occasion la tension alternative s'impose-t-elle définitivement ? \hfill \textcolor{red}{(1 pt)}

{\small \textcolor{blue}{\textbf{Réponse :} Le courant alternatif s'impose définitivement lors de l'Exposition universelle de Chicago en 1893, où Westinghouse et Tesla illuminent l'exposition avec leur système, démontrant ainsi sa supériorité. Cette victoire est confirmée avec le contrat pour équiper les chutes du Niagara d'une centrale hydroélectrique en courant alternatif.}}

\item Dans quelle grande autre aventure scientifique N Tesla se lance-t-il ? \hfill \textcolor{red}{(0,5 pt)}

{\small \textcolor{blue}{\textbf{Réponse :} Tesla se lance dans la transmission d'énergie et de signaux sans fil (radio), ainsi que dans de nombreuses recherches sur les hautes fréquences et la résonance électrique.}}

\item Quel est le nom et la nationalité de son concurrent ? \hfill \textcolor{red}{(0,5 pt)}

{\small \textcolor{blue}{\textbf{Réponse :} Son concurrent est Guglielmo Marconi, de nationalité italienne. Marconi est souvent crédité de l'invention de la radio, bien que Tesla ait déposé des brevets antérieurs.}}

\item Combien de brevets industriels N Tesla a-t-il déposé ? \hfill \textcolor{red}{(0,5 pt)}

{\small \textcolor{blue}{\textbf{Réponse :} Nicola Tesla a déposé environ 300 brevets au cours de sa vie (le nombre exact varie selon les sources entre 278 et plus de 300 brevets, dans différents pays).}}

\end{enumerate}

\vspace{2mm}
\begin{center}
\large{\textcolor{red}{\textbf{TOTAL : / 10 points}}}
\end{center}

\end{document}